\section{Threat Model}\label{sec:threat}
\paragraph{Assets and adversary}
The private target is the user's requested task, here one of six support actions. The defender controls the Agent and its controller, and may run relays, but does not control third-party tool services. We consider three observers, distinguished by their rights. A \emph{transit observer} sees connection destinations, directions, ciphertext lengths and times on the links it watches. A \emph{recipient} additionally reads the requests addressed to it, including arguments. An \emph{active auditor} also chooses probes. Identifying a requested action does not establish whether it was authorized; that needs the principal and policy context. Table~\ref{tab:deployment-access} maps these rights to common deployments and to the existing evidence for each view.

\begin{table*}[t]\centering\small
\caption{Deployment and access map. The configurations are analyst-defined; cited settings test the listed view, not every component of the deployment. A gateway's plaintext access requires TLS termination.}\label{tab:deployment-access}
\setlength{\tabcolsep}{4pt}
\begin{tabularx}{\textwidth}{@{}P{.18\textwidth}P{.25\textwidth}P{.28\textwidth}Y@{}}\toprule
Deployment & Observer and permitted view & Existing evidence for that view & Protection consequence \\\midrule
Enterprise egress & A transit monitor sees routed model/tool flows, sizes, times and any exposed destination association. & WRA: supplied research-domain traces; ANT and cloud benchmark: model/combined traffic~\cite{jeong2026wra,li2026ant,pouryousef2026traffic}. & Specify which services cross the monitored egress. A union of flows is available only when the monitor can observe them. \\\addlinespace[3pt]
Shared gateway or relay & Transit sees multiplexed gateway output. A terminating gateway can additionally read addressed requests. & Live OHTTP: passive-link inference under stated relay trust~\cite{dangol2026workflow}; PrivacyTrace: semantic recipient/gateway access~\cite{duan2026privacytrace}. & Separate relay association protection from content trust. Connection sharing changes the observation unit. \\\addlinespace[3pt]
Local model, remote tools & An external observer sees remote tool traffic; the local model call has no corresponding external network leg. & WRA's destination-recovery mechanism is relevant to tool egress~\cite{jeong2026wra}; whole-deployment protection is untested in the mapped comparison. & Evaluate tool and downstream calls. Do not import a remote-model attack score into an absent network leg. \\\addlinespace[3pt]
Cloud model, remote tools & An authorized vantage may see either or both encrypted legs. A contacted tool sees its own arguments. & ANT/cloud benchmark test model traffic~\cite{li2026ant,pouryousef2026traffic}; PrivacyTrace tests recipient-visible representations~\cite{duan2026privacytrace}. & Check permitted model and tool views after the intervention; recipients and passive monitors need different protections. \\\bottomrule
\end{tabularx}
\end{table*}


\paragraph{When cover must execute}
Transport cover ends before a business tool, at an endpoint the defender controls~\cite{apthorpe2019stp}. That endpoint must sit after the observer, or the observer sees it discard the cover. A defender can move the endpoint, for example to a relay outside a monitored egress, but the segment from the last defender-controlled hop to a third-party service remains. If the adversary observes that segment, or is the service, every request it sees must be real (Table~\ref{tab:cover-necessity}). Executed cover then follows from the deployment rather than from a design preference.

\begin{table}[t]\centering\footnotesize
\caption{When cover must execute. Transport cover needs a defender-controlled terminator after the observer and before the business service. A relay can move the terminator but not past the final hop to a third-party service; an observer of that hop, or the service itself, sees only real calls.}\label{tab:cover-necessity}
\setlength{\tabcolsep}{3pt}
\begin{tabularx}{\columnwidth}{@{}P{.27\columnwidth}P{.2\columnwidth}Y@{}}\toprule
Observer $\to$ service & Terminator after observer? & Cover that can work \\\midrule
Egress $\to$ own service & Yes (own tunnel) & Transport cover suffices~\cite{apthorpe2019stp} \\\addlinespace[2pt]
Egress $\to$ third-party SaaS or remote MCP & Only via a relay outside the egress & Transport cover against the egress alone; executed cover if the relay's upstream hop is observed \\\addlinespace[2pt]
Recipient (tool provider) & No & Executed cover whose arguments must also be plausible to the recipient \\\addlinespace[2pt]
TLS-terminating gateway & After gateway only & Executed cover, or trust the gateway with content \\\bottomrule
\end{tabularx}\end{table}


\paragraph{Observer in this paper}
Our measurements use a transit observer on the final hop to seven tool services that stand for third-party tools, with no defender relay in between. It receives destinations, ciphertext bytes and TLS record completion times. Section~\ref{sec:dispatch} additionally supplies episode and call windows, an evaluator privilege that Section~\ref{sec:noboundary} withdraws. The recipient view, where decoy arguments are also visible, is part of the threat model but not of these measurements (Section~\ref{sec:limits}). The adversary may train on traffic generated under the same cover policy; the attack suite therefore includes classifiers and cover-removal models fitted on protected training executions, alongside attacks frozen on unprotected traffic.
